% GENERATED FILE. Edit canonical lessons, then run npm run generate:books.
% canonical-source-hash: fnv1a64:56c0b35b03ed26aa
% canonical-lessons: ZH-C153-bi, ZH-W153-bi-compare, ZH-C153-bi-grammar, ZH-C153-geng, ZH-C153-gen, ZH-C153-name, ZH-C153-practice

\chapter{Older Than Me}
\label{ch:zh-older-than-me}
\hlchaptermodality{153}

\begin{chapteropening}
\textbf{By the end of this chapter:} \emph{I can compare two people or things.}

Comparing two people or things with \zh{比}, saying by how much, even more, the same as, and not as much as, built on one new written character: \zh{比}.
\end{chapteropening}

\section[\texorpdfstring{\zh{比} (bǐ)}{比 (bǐ)}]{\zh{比} (bǐ) — than; compared with}

\label{lesson:ZH-C153-bi}



\begin{quote}
A new chapter, and it opens on the word that sets two things side by side.

\begin{itemize}
\raggedright
  \item \emph{Recall:} say \zh{他说汉语说得很快} — \emph{tā shuō Hànyǔ shuō de hěn kuài} — \textbf{R1}, the chapter you have finished
  \item \emph{Recall:} say \emph{qīngchu}, clear — \textbf{R1}
  \item \emph{Recall:} say \emph{yíyàng}, the same — \textbf{R3}, durable
  \item \emph{Recall:} say \emph{bǐ}, a pen — \textbf{R2}, a few chapters back
\end{itemize}
\end{quote}

\subsection*{You'll want to know: \zh{比}}
\textbf{\zh{比}} — \emph{bǐ} — \textquotedblleft{}than; compared with\textquotedblright{}.

\begin{quote}
\textbf{\zh{我比他大}} \emph{wǒ bǐ tā dà} — I'm older than him.
\end{quote}

The next lesson puts the character in your hand. It sounds like \emph{bǐ}, a pen, and it is a different word.

\begin{cousinweb}
\textbf{\zh{比}} (\emph{bǐ}) is inherited from Middle Chinese \emph{pjijX} (Guangyun, in Baxter's transcription), glossed 'to compare; side by side'. It is the bǐ of \emph{bǐsài}, a match. The \emph{bǐ} of \emph{qiānbǐ}, a pencil, is another word with the same sound.
\end{cousinweb}

\begin{grammarlens}[title={What you've built}]
\zh{比} sits between the two things compared, and the adjective comes last: \zh{我}, \zh{比}, \zh{他}, \zh{大}.

\zh{大} for people is old: \zh{我比他大} is I'm older than him. \zh{比} itself is the bǐ of \emph{bǐsài}, a match: two sides set against each other.
\end{grammarlens}

\subsection*{Your turn}
\begin{itemize}
\raggedright
  \item \emph{Say it:} \emph{bǐ}
  \item \emph{Say it:} \emph{bǐ}, once more
  \item \emph{Say it:} \zh{我比他大} — \emph{wǒ bǐ tā dà} — I'm older than him.
  \item \emph{Say it:} \zh{他比我大} — \emph{tā bǐ wǒ dà} — He's older than me.
  \item \emph{Recall:} say \emph{bǐ}, a pen, then \zh{比}, than: one sound, two words
\end{itemize}

\begin{tcolorbox}[breakable,colback=teal!4,colframe=teal!35!black,title={Before you move on}]
What does \zh{比} mean? (\textquotedblleft{}Than; compared with\textquotedblright{}.) Where does the adjective go? (\textbf{Last: \zh{我比他大}}.)
\end{tcolorbox}
\section[\texorpdfstring{\zh{比} (bǐ)}{比 (bǐ)}]{\zh{比} — two halves side by side, the left one first}

\label{lesson:ZH-W153-bi-compare}



\begin{quote}
The word first, then the character inside it.

\begin{itemize}
\raggedright
  \item \emph{Recall:} say \emph{bǐ} — \textbf{R1}, one lesson back, before it decays
\end{itemize}

\begin{quote}
\zh{能} \zh{七}
\end{quote}

\emph{Read it:} each one aloud, with nothing said first — \emph{néng}, \emph{qī}
\end{quote}

\subsection*{Script — two halves side by side, the left one first}
\hlblockfigure{\detokenize{figures/ZH-W153-bi-compare-filmstrip.pdf}}{How \zh{比} is written, stroke by stroke}

\begin{quote}
\zh{比}
\end{quote}

\textbf{Two halves, side by side, the left one first.} The right half you write already, twice over, on the right of \zh{能}.

\textbf{The left half, two strokes}: a short bar, then a vertical that runs down and turns up to the right at its foot.

\textbf{The right half, two strokes}, as on the right of \zh{能}: a short left-falling stroke, then a vertical that bends right along the foot and hooks up, the same movement that ends \zh{七}.

2 + 2 = four.

\subsection*{Your turn}
\begin{itemize}
\raggedright
  \item \emph{Write it:} \zh{比}, all four, with the model in front of you
  \item \emph{Cover:} write it again from memory, and count the pieces, not the strokes
\end{itemize}

\begin{tcolorbox}[breakable,colback=teal!4,colframe=teal!35!black,title={Before you move on}]
How many strokes in \zh{比}? (\textbf{Four}.) Which half comes first? (\textbf{The left}.) Where have you written the right half? (\textbf{On the right of \zh{能}}.)

Source: \href{https://raw.githubusercontent.com/chanind/hanzi-writer-data/68d10a4b21150cae5e1ebbd223eed289cf32d90c/data/\%E6\%AF\%94.json}{Hanzi Writer Data \zh{比}, pinned stroke paths and medians}.
\end{tcolorbox}
\section[\texorpdfstring{\zh{他比我大两岁} (tā bǐ wǒ dà liǎng suì)}{他比我大两岁 (tā bǐ wǒ dà liǎng suì)}]{\zh{他比我大两岁} — comparing with \zh{比}}

\label{lesson:ZH-C153-bi-grammar}



\begin{quote}
Two people, one comparison, and then by how much.

\begin{itemize}
\raggedright
  \item \emph{Read it:} \zh{我比他大}, aloud — \emph{wǒ bǐ tā dà}
  \item \emph{Recall:} say \zh{二十岁} — \emph{èrshí suì}, twenty years old
  \item \emph{Recall:} say \emph{zuì}, most
\end{itemize}
\end{quote}

\begin{grammarlens}[title={Grammar Lens: \zh{比}, then the adjective, then by how much}]
\begin{quote}
\textbf{A, \zh{比}, B, then the adjective: A is more so than B. How much more goes after the adjective, at the very end.}
\end{quote}

\begin{quote}
\textbf{\zh{我比他大}} \emph{wǒ bǐ tā dà} — I'm older than him. \textbf{\zh{他比我大两岁}} \emph{tā bǐ wǒ dà liǎng suì} — He's two years older than me. \textbf{\zh{这个比那个大一点儿}} \emph{zhège bǐ nàge dà yìdiǎnr} — This one is a bit bigger than that one.
\end{quote}

There is no \zh{很} in these. \zh{很} describes one thing on its own; a comparison sets two against each other. \zh{他很大} describes him; \zh{他比我大} measures him against me.

The other way round, not as much as, comes with \zh{没有}:

\begin{quote}
\textbf{\zh{我没有他那么大}} \emph{wǒ méiyǒu tā nàme dà} — I'm not as old as he is.
\end{quote}

And \emph{zuì}, most, needs no second thing at all: \emph{tā zuì dà}.
\end{grammarlens}

\subsection*{Your turn}
Say these aloud:

\begin{itemize}
\raggedright
  \item \zh{他比我大两岁} — \emph{tā bǐ wǒ dà liǎng suì} — He's two years older than me.
  \item \zh{这个比那个大一点儿} — \emph{zhège bǐ nàge dà yìdiǎnr} — This one is a bit bigger than that one.
  \item \zh{我没有他那么大} — \emph{wǒ méiyǒu tā nàme dà} — I'm not as old as he is.
  \item \emph{jīntiān bǐ míngtiān rè} — Today is hotter than tomorrow.
\end{itemize}

\begin{tcolorbox}[breakable,colback=teal!4,colframe=teal!35!black,title={Before you move on}]
Where does \zh{比} go? (\textbf{Between the two things}.) Where does by how much go? (\textbf{At the very end: \zh{大两岁}}.) Does \zh{很} go in? (\textbf{No}.)
\end{tcolorbox}
\section[gèng]{gèng — even more}

\label{lesson:ZH-C153-geng}



\begin{quote}
Before the new one: read the last lesson's sentence, then a word for tall.

\begin{itemize}
\raggedright
  \item \emph{Read it:} \zh{他比我大两岁}, aloud — \emph{tā bǐ wǒ dà liǎng suì}
  \item \emph{Recall:} say \emph{gāo}, tall — \textbf{R3}, durable
\end{itemize}
\end{quote}

\subsection*{You'll want to know: gèng}
\emph{gèng} — \textquotedblleft{}even more; still more\textquotedblright{}.

\begin{quote}
\emph{tā gèng gāo} — He's even taller.
\end{quote}

\begin{quote}
\emph{tā bǐ wǒ gèng gāo} — He's even taller than me.
\end{quote}

Its character comes later; for now it is a word for the ear.

\begin{cousinweb}
\emph{gèng} is inherited from Middle Chinese \emph{kaengH} (Guangyun, in Baxter's transcription), glossed 'to change; to alter'. 'Even more' is a later sense, assumed here, not quoted.
\end{cousinweb}

\begin{grammarlens}[title={What you've built}]
\emph{gèng} goes straight before the adjective: \emph{tā gèng gāo}.

With \zh{比} it sharpens the comparison: both are tall, and he is taller still.
\end{grammarlens}

\subsection*{Your turn}
\begin{itemize}
\raggedright
  \item \emph{Say it:} \emph{gèng}
  \item \emph{Say it:} \emph{gèng}, once more
  \item \emph{Say it:} \emph{tā gèng gāo} — He's even taller.
  \item \emph{Say it:} \emph{tā bǐ wǒ gèng gāo} — He's even taller than me.
  \item \emph{Recall:} say \emph{gāo}, then \emph{gèng gāo}
\end{itemize}

\begin{tcolorbox}[breakable,colback=teal!4,colframe=teal!35!black,title={Before you move on}]
What does \emph{gèng} mean? (\textquotedblleft{}Even more\textquotedblright{}.) Where does it go? (\textbf{Straight before the adjective}.)
\end{tcolorbox}
\section[gēn]{gēn — with; the same as}

\label{lesson:ZH-C153-gen}



\begin{quote}
Before the new one: last lesson's word, then the word for the same.

\begin{itemize}
\raggedright
  \item \emph{Recall:} say \emph{gèng}, even more — \textbf{R1}, one lesson back
  \item \emph{Recall:} say \emph{yíyàng}, the same — \textbf{R3}, durable
\end{itemize}
\end{quote}

\subsection*{You'll want to know: gēn}
\emph{gēn} — \textquotedblleft{}with; and\textquotedblright{}; with \emph{yíyàng} after it, \textquotedblleft{}the same as\textquotedblright{}.

\begin{quote}
\emph{wǒ gēn tā yíyàng gāo} — I'm as tall as he is.
\end{quote}

\begin{quote}
\emph{wǒ gēn māma qù} — I'm going with Mum.
\end{quote}

Its character comes later; for now it is a word for the ear.

\begin{cousinweb}
\emph{gēn} is inherited from Middle Chinese \emph{kon} (Guangyun, in Baxter's transcription), glossed 'the heel'. How the heel became 'with', following at someone's heels, is a guess the gloss does not make, so the history is not settled.
\end{cousinweb}

\begin{grammarlens}[title={What you've built}]
\emph{gēn} … \emph{yíyàng}, then an adjective, says two things are the same: \emph{wǒ gēn tā yíyàng gāo}.

\zh{比} says one is more; \emph{gēn} … \emph{yíyàng} says neither is.
\end{grammarlens}

\subsection*{Your turn}
\begin{itemize}
\raggedright
  \item \emph{Say it:} \emph{gēn}
  \item \emph{Say it:} \emph{gēn}, once more
  \item \emph{Say it:} \emph{wǒ gēn tā yíyàng gāo} — I'm as tall as he is.
  \item \emph{Say it:} \emph{wǒ gēn māma qù} — I'm going with Mum.
  \item \emph{Recall:} say \emph{yíyàng}, then \emph{gēn tā yíyàng}
\end{itemize}

\begin{tcolorbox}[breakable,colback=teal!4,colframe=teal!35!black,title={Before you move on}]
What does \emph{gēn} mean? (\textquotedblleft{}With; and\textquotedblright{}.) And \emph{gēn} … \emph{yíyàng}? (\textbf{The same as}.)
\end{tcolorbox}
\section[\texorpdfstring{\zh{那么} (nàme)}{那么 (nàme)}]{\zh{那么} (nàme) — so; that much}

\label{lesson:ZH-C153-name}



\begin{quote}
Before the new one: last lesson's word, then the two words this one is made of.

\begin{itemize}
\raggedright
  \item \emph{Recall:} say \emph{gēn}, with — \textbf{R1}, one lesson back
  \item \emph{Recall:} say \zh{那} — \emph{nà}, that — \textbf{R4}, recognition at distance
  \item \emph{Recall:} say \zh{什么} — \emph{shénme}, what — \textbf{R4}
\end{itemize}
\end{quote}

\subsection*{You'll want to know: \zh{那么}}
\textbf{\zh{那么}} — \emph{nàme} — \textquotedblleft{}so; that (much)\textquotedblright{}.

\begin{quote}
\textbf{\zh{我没有他那么大}} \emph{wǒ méiyǒu tā nàme dà} — I'm not as old as he is.
\end{quote}

\begin{quote}
\textbf{\zh{那么大}} \emph{nàme dà} — that big
\end{quote}

\begin{cousinweb}
\textbf{\zh{那么}} (\emph{nàme}) is built from \zh{那} 'that' and \zh{么}, the \zh{么} of \zh{什么}. The Guangyun records \zh{那} as \emph{na} (Guangyun, in Baxter's transcription), glossed 'how; what'; how it came to mean 'that', and where the ending \zh{么} comes from, are not settled.
\end{cousinweb}

\begin{grammarlens}[title={What you've built}]
\zh{没有} … \zh{那么}, with an adjective after it, says not as much as: \zh{我没有他那么大}.

It is the other side of \zh{比}: \zh{他比我大}, so \zh{我没有他那么大}.
\end{grammarlens}

\subsection*{Your turn}
\begin{itemize}
\raggedright
  \item \emph{Say it:} \emph{nàme}
  \item \emph{Say it:} \emph{nàme}, once more
  \item \emph{Say it:} \zh{我没有他那么大} — \emph{wǒ méiyǒu tā nàme dà} — I'm not as old as he is.
  \item \emph{Say it:} \zh{他比我大} — \emph{tā bǐ wǒ dà} — He's older than me.
  \item \emph{Recall:} say \zh{那} and \zh{什么}, then \zh{那么}
\end{itemize}

\begin{tcolorbox}[breakable,colback=teal!4,colframe=teal!35!black,title={Before you move on}]
What does \zh{那么} mean? (\textquotedblleft{}So; that much\textquotedblright{}.) Say you are not as old as him. (\textbf{\zh{我没有他那么大}}.)
\end{tcolorbox}
\section[Practice]{\zh{我比他大} — \zh{我没有他那么大}}

\label{lesson:ZH-C153-practice}



\begin{quote}
The chapter closes here, with two people comparing a family.

\begin{itemize}
\raggedright
  \item \emph{Recall:} say \zh{那么} — \emph{nàme}, so — \textbf{R1}, one lesson back
  \item \emph{Recall:} say \emph{gēn} … \emph{yíyàng}, the same as — \textbf{R1}
  \item \emph{Recall:} say \emph{gèng}, even more — \textbf{R2}
\end{itemize}
\end{quote}

\subsection*{The exchange}
\begin{quote}
\textbf{A:} \zh{你多大}
\end{quote}

\begin{quote}
\emph{nǐ duō dà} — How old are you?
\end{quote}

\begin{quote}
\textbf{B:} \zh{我二十岁}
\end{quote}

\begin{quote}
\emph{wǒ èrshí suì} — I'm twenty.
\end{quote}

\begin{quote}
\textbf{A:} \zh{他呢}
\end{quote}

\begin{quote}
\emph{tā ne} — And him?
\end{quote}

\begin{quote}
\textbf{B:} \zh{他比我大两岁}
\end{quote}

\begin{quote}
\emph{tā bǐ wǒ dà liǎng suì} — He's two years older than me.
\end{quote}

\begin{quote}
\textbf{A:} \emph{nǐ bǐ tā gāo ma} — Are you taller than him?
\end{quote}

\begin{quote}
\textbf{B:} \emph{wǒ méiyǒu tā nàme gāo} — I'm not as tall as he is.
\end{quote}

\begin{quote}
\textbf{A:} \zh{你妈妈呢}
\end{quote}

\begin{quote}
\emph{nǐ māma ne} — And your mum?
\end{quote}

\begin{quote}
\textbf{B:} \emph{tā gēn wǒ yíyàng gāo} — She's as tall as me.
\end{quote}

\begin{quote}
\textbf{A:} \emph{wǒ bǐ tā dà, tā gèng gāo} — I'm older than him, and he's taller still.
\end{quote}

\zh{比} puts the difference last: \zh{大两岁}. \zh{没有} … \zh{那么} turns it round. \emph{gēn} … \emph{yíyàng} says neither is more.

\begin{tcolorbox}[breakable,skin=enhanced,colback=orange!6,colframe=orange!45!black,boxrule=0.5pt,arc=1mm,left=6pt,right=6pt,top=4pt,bottom=4pt,fonttitle=\bfseries,title={Read this}]
\emph{Read it:} the exchange through once without stopping, then answer the question below
\end{tcolorbox}

\begin{grammarlens}[title={What you've built — more, less and the same}]
Three ways to set two things side by side:

\begin{itemize}
\raggedright
  \item \textbf{\zh{他比我大两岁}} \emph{tā bǐ wǒ dà liǎng suì} — he's two years older than me.
  \item \textbf{\zh{我没有他那么大}} \emph{wǒ méiyǒu tā nàme dà} — I'm not as old as he is.
  \item \emph{wǒ gēn tā yíyàng gāo} — I'm as tall as he is.
\end{itemize}

\textbf{Listen.} \emph{Nǐ bǐ nǐ māma gāo ma?}

\textbf{Speak.} Compare yourself with somebody you know: older or younger, taller or not as tall.

\emph{Write it:} \zh{比} — the left half first, then the right half of \zh{能}
\end{grammarlens}

\subsection*{Your turn}
\begin{itemize}
\raggedright
  \item \emph{Say it:} \zh{他比我大两岁} — \emph{tā bǐ wǒ dà liǎng suì} — He's two years older than me.
  \item \emph{Say it:} \zh{我没有他那么大} — \emph{wǒ méiyǒu tā nàme dà} — I'm not as old as he is.
  \item \emph{Say it:} the whole exchange, both parts
  \item \emph{Write it:} \zh{我比他大}
\end{itemize}

\begin{tcolorbox}[breakable,colback=teal!4,colframe=teal!35!black,title={Before you move on}]
Say you are older than him. (\textbf{\zh{我比他大}}.) Say you are not as old as him. (\textbf{\zh{我没有他那么大}}.) Say you are as tall as he is. (\emph{wǒ gēn tā yíyàng gāo}.)
\end{tcolorbox}
